\paragraph*{04.05.01.05 Konzept für die Projektdurchführung entwerfen, überwachen und anpassen}
\kcilabel{para:04.05.01.05_KonzeptFürDieProjektdurchführungEntwerfenÜberwachenUndAnpassen}{04.05.01.05}
\label{para:04.05.01.05_KonzeptFürDieProjektdurchführungEntwerfenÜberwachenUndAnpassen}

% Task
\par Nach der Vorstandentscheidung des \gls{adac}, die \gls{adn} als \gls{bk} in das \gls{f2h}-Projekt zu integrieren, war eine Anpassung des Projektdurchführungskonzepts erforderlich. Ziel war es, die neuen \gls{adn}-Schnittstellen zu berücksichtigen und die Zusammenarbeit aller Parteien auf einer einheitlichen Systembasis zu sichern.

% Action
\par Die Anbindung der \gls{adn} an die \gls{erp}-Systeme des \gls{adac} war neuartig; etablierte Prozesse oder Best Practices fehlten. Deshalb entwickelte ich ein neues Durchführungs-Konzept, das die speziellen Anforderungen und Herausforderungen der \gls{adn}-Integration umfasst. Die Entscheidung fiel auf die Nutzung der für das Hauptprojekt \gls{r2h} geplanten \gls{aif}-Schnittstellentechnik für die \gls{adn}-Schnittstellen.

\par Zur Risikominimierung arbeitete das \gls{f2h}-Projekt mit kürzeren Iterationen und einem agilen Projektmanagementtool. Ziel war die frühzeitige Implementierung und Prüfung der \gls{adn}-Schnittstellen. Regelmäßige Abstimmungen stellten sicher, dass alle Partner informiert waren und Probleme frühzeitig adressiert werden konnten.

% Result
\par Die Anpassung des Durchführungs-Konzepts ermöglichte die erfolgreiche Berücksichtigung der \gls{adn}-Integration. Das Projektmanagementtool erlaubte kontinuierliche Überwachung der Schnittstellenentwicklung und gezielte Anpassungen. Dieser Zuschnitt findet auch im Hauptprojekt \gls{r2h} für die Welle 1.2 ff. Anwendung.

% Reflect
\par Die Integration der \gls{adn} war technisch, organisatorisch und menschlich herausfordernd, weil die \gls{adn} eine \gls{s4}-Anbindung nicht bevorzugte. Aufgrund eines schwachen Mandats konnte ich die \gls{adn}-Anbindung nicht durchsetzen.

\par Dennoch sicherte das angepasste Konzept eine erfolgreiche Integration der \gls{adn} ohne größere Verzögerungen oder Probleme. Wünschenswert wäre gewesen, die \gls{adn}-Anbindung von Beginn an als festen Projektbestandteil zu berücksichtigen, um eine nahtlosere Integration und frühere Problemlösungen zu ermöglichen.